\chapter{The Messenger}

\lettrine{M}{ademoiselle} de Montalais was right; the young cavalier was goodly to look upon.

\zz
He was a young man of from twenty-four to twenty-five years of age, tall and slender, wearing gracefully the picturesque military costume of the period. His large boots contained a foot which Mademoiselle de Montalais might not have disowned if she had been transformed into a man. With one of his delicate but nervous hands he checked his horse in the middle of the court, and with the other raised his hat, whose long plumes shaded his at once serious and ingenuous countenance.

The guards, roused by the steps of the horse, awoke, and were on foot in a minute. The young man waited till one of them was close to his saddle-bow: then, stooping towards him, in a clear, distinct voice, which was perfectly audible at the window where the two girls were concealed, <A message for his royal highness,> he said.

<Ah, ah!> cried the soldier. <Officer, a messenger!>

But this brave guard knew very well that no officer would appear, seeing that the only one who could have appeared dwelt at the other side of the castle, in an apartment looking into the gardens. So he hastened to add: <The officer, monsieur, is on his rounds; but, in his absence, M.~de Saint-Rémy, the maître d'hôtel, shall be informed.>

<M.~de Saint-Rémy?> repeated the cavalier, slightly blushing.

<Do you know him?>

<Why, yes; but request him, if you please, that my visit be announced to his royal highness as soon as possible.>

<It appears to be pressing,> said the guard, as if speaking to himself, but really in the hope of obtaining an answer.

The messenger made an affirmative sign with his head.

<In that case,> said the guard, <I will go and seek the maître d'hôtel myself.>

The young man, in the meantime, dismounted; and whilst the others were making their remarks upon the fine horse the cavalier rode, the soldier returned.

<Your pardon, young gentleman; but your name, if you please?>

<The Vicomte de Bragelonne, on the part of his highness M.~le Prince de Condé.>

The soldier made a profound bow, and, as if the name of the conqueror of Rocroi and Lens had given him wings, he stepped lightly up the steps leading to the ante-chamber.

A. de Bragelonne had not had time to fasten his horse to the iron bars of the perron, when M.~de Saint-Rémy came running, out of breath, supporting his capacious body with one hand, whilst with the other he cut the air as a fisherman cleaves the waves with his oar.

<Ah, Monsieur le Vicomte! You at Blois!> cried he. <Well, that is a wonder. Good-day to you—good-day, Monsieur Raoul.>

<I offer you a thousand respects, M.~de Saint-Rémy.>

<How Madame de la Vall—I mean, how delighted Madame de Saint-Rémy will be to see you! But come in. His royal highness is at breakfast—must he be interrupted? Is the matter serious?>

<Yes, and no, Monsieur de Saint-Rémy. A moment's delay, however, would be disagreeable to his royal highness.>

<If that is the case, we will force the consigne, Monsieur le Vicomte. Come in. Besides, Monsieur is in an excellent humour to-day. And then you bring news, do you not?>

<Great news, Monsieur de Saint-Rémy.>

<And good, I presume?>

<Excellent.>

<Come quickly, come quickly then!> cried the worthy man, putting his dress to rights as he went along.

Raoul followed him, hat in hand, and a little disconcerted at the noise made by his spurs in these immense salons.

As soon as he had disappeared in the interior of the palace, the window of the court was repeopled, and an animated whispering betrayed the emotion of the two girls. They soon appeared to have formed a resolution, for one of the two faces disappeared from the window. This was the brunette; the other remained behind the balcony, concealed by the flowers, watching attentively through the branches the perron by which M.~de Bragelonne had entered the castle.

In the meantime the object of so much laudable curiosity continued his route, following the steps of the maître d'hôtel. The noise of quick steps, an odour of wine and viands, a clinking of crystal and plates, warned them that they were coming to the end of their course.

The pages, valets and officers, assembled in the office which led up to the refectory, welcomed the newcomer with the proverbial politeness of the country; some of them were acquainted with Raoul, and all knew that he came from Paris. It might be said that his arrival for a moment suspended the service. In fact, a page, who was pouring out wine for his royal highness, on hearing the jingling of spurs in the next chamber, turned round like a child, without perceiving that he was continuing to pour out, not into the glass, but upon the tablecloth.

Madame, who was not so preoccupied as her glorious spouse was, remarked this distraction of the page.

<Well?> exclaimed she.

<Well!> repeated Monsieur; <what is going on then?>

A. de Saint-Rémy, who had just introduced his head through the doorway, took advantage of the moment.

<Why am I to be disturbed?> said Gaston, helping himself to a thick slice of one of the largest salmon that had ever ascended the Loire to be captured between Paimbœuf and Saint-Nazaire.

<There is a messenger from Paris. Oh! but after monseigneur has breakfasted will do; there is plenty of time.>

<From Paris!> cried the prince, letting his fork fall. <A messenger from Paris, do you say? And on whose part does this messenger come?>

<On the part of M.~le Prince,> said the maître d'hôtel promptly.

Every one knows that the Prince de Condé was so called.

<A messenger from M.~le Prince!> said Gaston, with an inquietude that escaped none of the assistants, and consequently redoubled the general curiosity.

Monsieur, perhaps, fancied himself brought back again to the happy times when the opening of a door gave him an emotion, in which every letter might contain a state secret,—in which every message was connected with a dark and complicated intrigue. Perhaps, likewise, that great name of M.~le Prince expanded itself, beneath the roofs of Blois, to the proportions of a phantom.

Monsieur pushed away his plate.

<Shall I tell the envoy to wait?> asked M.~de Saint-Rémy.

A glance from Madame emboldened Gaston, who replied: <No, no! let him come in at once, on the contrary. À propos, who is he?>

<A gentleman of this country, M.~le Vicomte de Bragelonne.>

<Ah, very well! Introduce him, Saint-Rémy—introduce him.>

And when he had let fall these words, with his accustomed gravity, Monsieur turned his eyes, in a certain manner, upon the people of his suite, so that all, pages, officers, and equerries, quitted the service, knives and goblets, and made towards the second chamber door a retreat as rapid as it was disorderly.

This little army had dispersed in two files when Raoul de Bragelonne, preceded by M.~de Saint-Rémy, entered the refectory.

The short interval of solitude which this retreat had left him, permitted Monsieur the time to assume a diplomatic countenance. He did not turn round, but waited till the maître d'hôtel should bring the messenger face to face with him.

Raoul stopped even with the lower end of the table, so as to be exactly between Monsieur and Madame. From this place he made a profound bow to Monsieur, and a very humble one to Madame; then, drawing himself up into military pose, he waited for Monsieur to address him.

On his part the prince waited till the doors were hermetically closed; he would not turn round to ascertain the fact, as that would have been derogatory to his dignity, but he listened with all his ears for the noise of the lock, which would promise him at least an appearance of secrecy.

The doors being closed, Monsieur raised his eyes towards the vicomte, and said, <It appears that you come from Paris, monsieur?>

<This minute, monseigneur.>

<How is the king?>

<His majesty is in perfect health, monseigneur.>

<And my sister-in-law?>

<Her majesty the queen-mother still suffers from the complaint in her chest, but for the last month she has been rather better.>

<Somebody told me you came on the part of M.~le Prince. They must have been mistaken, surely?>

<No, monseigneur; M.~le Prince has charged me to convey this letter to your royal highness, and I am to wait for an answer to it.>

Raoul had been a little annoyed by this cold and cautious reception, and his voice insensibly sank to a low key.

The prince forgot that he was the cause of this apparent mystery, and his fears returned.

He received the letter from the Prince de Condé with a haggard look, unsealed it as he would have unsealed a suspicious packet, and in order to read it so that no one should remark the effects of it upon his countenance, he turned round.

Madame followed, with an anxiety almost equal to that of the prince, every manœver of her august husband.

Raoul, impassible, and a little disengaged by the attention of his hosts, looked from his place through the open window at the gardens and the statues which peopled them.

<Well!> cried Monsieur, all at once, with a cheerful smile; <here is an agreeable surprise, and a charming letter from M.~le Prince. Look, Madame!>

The table was too large to allow the arm of the prince to reach the hand of Madame; Raoul sprang forward to be their intermediary, and did it with so good a grace as to procure a flattering acknowledgement from the princess.

<You know the contents of this letter, no doubt?> said Gaston to Raoul.

<Yes, monseigneur; M.~le Prince at first gave me the message verbally, but upon reflection his highness took up his pen.>

<It is beautiful writing,> said Madame, <but I cannot read it.>

<Will you read it to Madame, M.~de Bragelonne?> said the duke.

<Yes; read it, if you please, monsieur.>

Raoul began to read, Monsieur giving again all his attention. The letter was conceived in these terms:

\begin{mail}{}{Monseigneur—}
The king is about to set out for the frontiers. You are aware the marriage of his majesty is concluded upon. The king has done me the honour to appoint me his maréchal-des-logis for this journey, and as I knew with what joy his majesty would pass a day at Blois, I venture to ask your royal highness's permission to mark the house you inhabit as our quarters. If, however, the suddenness of this request should create to your royal highness any embarrassment, I entreat you to say so by the messenger I send, a gentleman of my suite, M.~le Vicomte de Bragelonne. My itinerary will depend on your royal highness's determination, and instead of passing through Blois, we shall come through Vendôme or Romorantin. I venture to hope that your royal highness will be pleased with my arrangement, it being the expression of my boundless desire to make myself agreeable to you.
\end{mail}

<Nothing can be more gracious toward us,> said Madame, who had more than once consulted the looks of her husband during the reading of the letter. <The king here!> exclaimed she, in a rather louder tone than would have been necessary to preserve secrecy.

<Monsieur,> said his royal highness in his turn, <you will offer my thanks to M.~de Condé, and express to him my gratitude for the honour he has done me.> Raoul bowed.

<On what day will his majesty arrive?> continued the prince.

<The king, monseigneur, will in all probability arrive this evening.>

<But how, then, could he have known my reply if it had been in the negative?>

<I was desired, monseigneur, to return in all haste to Beaugency, to give counter-orders to the courier, who was himself to go back immediately with counter-orders to M.~le Prince.>

<His majesty is at Orléans, then?>

<Much nearer, monseigneur; his majesty must by this time have arrived at Meung.>

<Does the court accompany him?>

<Yes, monseigneur.>

<À propos, I forgot to ask you after M.~le Cardinal.>

<His eminence appears to enjoy good health, monseigneur.>

<His nieces accompany him, no doubt?>

<No, monseigneur; his eminence has ordered the Mesdemoiselles de Mancini to set out for Brouage. They will follow the left bank of the Loire, while the court will come by the right.>

<What! Mademoiselle Mary de Mancini quit the court in that manner?> asked Monsieur, his reserve beginning to diminish.

<Mademoiselle Mary de Mancini in particular,> replied Raoul discreetly.

A fugitive smile, an imperceptible vestige of his ancient spirit of intrigue, shot across the pale face of the prince.

<Thanks, M.~de Bragelonne,> then said Monsieur. <You would, perhaps, not be willing to carry M.~le Prince the commission with which I would charge you, and that is, that his messenger has been very agreeable to me; but I will tell him so myself.>

Raoul bowed his thanks to Monsieur for the honour he had done him.

Monsieur made a sign to Madame, who struck a bell which was placed at her right hand; M.~de Saint-Rémy entered, and the room was soon filled with people.

<Messieurs,> said the prince, <his majesty is about to pay me the honour of passing a day at Blois; I depend on the king, my nephew, not having to repent of the favour he does my house.>

<Vive le Roi!> cried all the officers of the household with frantic enthusiasm, and M.~de Saint-Rémy louder than the rest.

Gaston hung down his head with evident chagrin. He had all his life been obliged to hear, or rather to undergo, this cry of <Vive le Roi!> which passed over him. For a long time, being unaccustomed to hear it, his ear had had rest, and now a younger, more vivacious, and more brilliant royalty rose up before him, like a new and more painful provocation.

Madame perfectly understood the sufferings of that timid, gloomy heart; she rose from the table, Monsieur imitated her mechanically, and all the domestics, with a buzzing like that of several bee-hives, surrounded Raoul for the purpose of questioning him.

Madame saw this movement, and called M.~de Saint-Rémy.

<This is not the time for gossiping, but working,> said she, with the tone of an angry housekeeper.

A. de Saint-Rémy hastened to break the circle formed by the officers round Raoul, so that the latter was able to gain the ante-chamber.

<Care will be taken of that gentleman, I hope,> added Madame, addressing M.~de Saint-Rémy.

The worthy man immediately hastened after Raoul. <Madame desires refreshments to be offered to you,> said he; <and there is, besides, a lodging for you in the castle.>

<Thanks, M.~de Saint-Rémy,> replied Raoul; <but you know how anxious I must be to pay my duty to M.~le Comte, my father.>

<That is true, that is true, Monsieur Raoul; present him, at the same time, my humble respects, if you please.>

Raoul thus once more got rid of the old gentleman, and pursued his way. As he was passing under the porch, leading his horse by the bridle, a soft voice called him from the depths of an obscure path.

\begin{bwbigpic}
	[\basicscale] 
	{02obscure}
	[A soft voice called him from the depths of an obscure path]
\end{bwbigpic}

<Monsieur Raoul!> said the voice.

The young man turned round, surprised, and saw a dark complexioned girl, who, with a finger on her lip, held out her other hand to him. This young lady was an utter stranger. 